\section{Frozen-Bias Reference Ensemble}

The implemented method separates two tasks that are easy to conflate:
\begin{enumerate}
    \item adaptive AWH sampling, which is used to construct a reliable reference
    expanded ensemble,
    \item offline thermodynamic estimation, which is performed only after that
    reference ensemble has been frozen.
\end{enumerate}
This distinction is essential. The optimization gradients in the current code are
\emph{not} obtained by differentiating the online AWH update rule. Instead, the
adaptive AWH phase is used only to produce a stationary reference distribution
from which replay-based estimators can be evaluated cleanly.

Throughout this section, $i \in \Lambda_\ell$ denotes a generic lambda-state
index, $m \in \Lambda_\ell$ denotes a target lambda state to which the stored
frames are reweighted, and $n$ indexes stored frames from the frozen reference
trajectory.

Consider one thermodynamic leg $\ell$ with lambda states
$i \in \Lambda_\ell$, inverse temperature $\beta_\ell$, and, for NPT legs,
pressure $P_\ell$. The reduced potential of state $i$ at parameters $\theta$ is
\begin{equation}
    u_\ell(\mathbf{x}, V, i; \theta)
    =
    \beta_\ell U_\ell(\mathbf{x}, i; \theta)
    +
    \beta_\ell P_\ell V,
    \label{eq:reducedpotential}
\end{equation}
with the pressure-volume term omitted for legs that do not include a $pV$
contribution.

During the frozen stage, the AWH bias and target lambda weights are fixed. The
reference Gibbs log-weights used by the implementation are
\begin{equation}
    g_{\mathrm{ref},\ell}(i)
    =
    f_\ell(i)
    +
    \log \rho_\ell(i),
    \label{eq:gref}
\end{equation}
where $f_\ell(i)$ is the frozen AWH free-energy estimate and $\rho_\ell(i)$ is
the fixed target lambda distribution. The corresponding frozen extended-ensemble
reference density is
\begin{equation}
    p_{\mathrm{ref},\ell}(\mathbf{x}, V, i)
    =
    \frac{
        \exp \left[
            g_{\mathrm{ref},\ell}(i)
            -
            u_{\mathrm{ref},\ell}(\mathbf{x}, V, i)
        \right]
    }{
        \mathcal{Z}_{\mathrm{ref},\ell}
    },
    \label{eq:pref}
\end{equation}
where $u_{\mathrm{ref},\ell}$ denotes the reduced potential at the frozen
reference parameterization and
$\mathcal{Z}_{\mathrm{ref},\ell}$ is the corresponding normalizing constant.

For a stored frame $n$ with coordinates $\mathbf{x}_n$ and volume $V_n$, the code
evaluates the reference-mixture denominator
\begin{equation}
    \log Z_{\mathrm{mix},\ell}(n)
    =
    \log
    \sum_{i \in \Lambda_\ell}
    \exp \left[
        g_{\mathrm{ref},\ell}(i)
        -
        u_{\mathrm{ref},\ell}(\mathbf{x}_n, V_n, i)
    \right].
    \label{eq:logmix}
\end{equation}
Equation \eqref{eq:logmix} is the key MBAR/AWH consistency relation in the
current implementation. It is the object against which replayed frames are
reweighted both for readiness analysis and for optimization.

\section{MBAR Reweighting to an Arbitrary Lambda State}

Let $m \in \Lambda_\ell$ be a target lambda state at which we want to evaluate a
free energy or a free-energy gradient under the current parameter vector
$\theta$. Define the target-state partition function
\begin{equation}
    Q_{\ell,m}(\theta)
    =
    \int dV \int d\mathbf{x} \,
    \exp \left[
        -u_\ell(\mathbf{x}, V, m; \theta)
    \right].
    \label{eq:qtarget}
\end{equation}
The corresponding reduced free energy is
\begin{equation}
    F_{\ell,m}(\theta)
    =
    -\log Q_{\ell,m}(\theta).
    \label{eq:ftarget}
\end{equation}

To connect $Q_{\ell,m}$ to the frozen reference ensemble, multiply and divide the
integrand in Equation \eqref{eq:qtarget} by the reference-mixture density:
\begin{align}
    Q_{\ell,m}(\theta)
    &=
    \int dV \int d\mathbf{x} \,
    \mathcal{Z}_{\mathrm{ref},\ell}
    p_{\mathrm{ref},\ell}(\mathbf{x}, V)
    \exp \left[
        -u_\ell(\mathbf{x}, V, m; \theta)
        -
        \log Z_{\mathrm{mix},\ell}(\mathbf{x}, V)
    \right]
    \nonumber \\
    &=
    \mathcal{Z}_{\mathrm{ref},\ell}
    \,
    \mathbb{E}_{\mathrm{ref},\ell}
    \left[
        \exp \left(
            -u_\ell(\mathbf{x}, V, m; \theta)
            -
            \log Z_{\mathrm{mix},\ell}(\mathbf{x}, V)
        \right)
    \right].
    \label{eq:qtargetref}
\end{align}
Here
$\log Z_{\mathrm{mix},\ell}(\mathbf{x}, V)$ denotes the framewise quantity in
Equation \eqref{eq:logmix}, but written as a function of a generic configuration
instead of an indexed stored frame.

Taking minus the logarithm of Equation \eqref{eq:qtargetref} gives
\begin{equation}
    F_{\ell,m}(\theta)
    =
    -\log \mathcal{Z}_{\mathrm{ref},\ell}
    -
    \log
    \mathbb{E}_{\mathrm{ref},\ell}
    \left[
        \exp \left(
            -u_\ell(\mathbf{x}, V, m; \theta)
            -
            \log Z_{\mathrm{mix},\ell}(\mathbf{x}, V)
        \right)
    \right].
    \label{eq:freeref}
\end{equation}
The first term depends only on the frozen reference ensemble and is therefore a
constant shared by every target state $m$ of the same leg. The implementation
does not need that constant explicitly because only free-energy differences are
used. It is therefore convenient to define the additive-constant-free quantity
\begin{equation}
    \widetilde{F}_{\ell,m}(\theta)
    =
    -
    \log
    \mathbb{E}_{\mathrm{ref},\ell}
    \left[
        \exp \left(
            -u_\ell(\mathbf{x}, V, m; \theta)
            -
            \log Z_{\mathrm{mix},\ell}(\mathbf{x}, V)
        \right)
    \right],
    \label{eq:fendpointdef}
\end{equation}
so that
\begin{equation}
    F_{\ell,m}(\theta)
    =
    -\log \mathcal{Z}_{\mathrm{ref},\ell}
    +
    \widetilde{F}_{\ell,m}(\theta).
    \label{eq:fendpointconst}
\end{equation}
All endpoint differences within a leg depend only on
$\widetilde{F}_{\ell,m}(\theta)$ because the reference constant cancels.

\section{Discrete MBAR Weights Used by the Code}

Suppose the frozen reference trajectory for leg $\ell$ contains
$M_\ell$ stored frames. Replacing the expectation in
Equation \eqref{eq:fendpointdef} by the empirical average over those frames gives
\begin{equation}
    \widetilde{F}_{\ell,m}(\theta)
    \approx
    -
    \log
    \left[
        \frac{1}{M_\ell}
        \sum_{n=1}^{M_\ell}
        a_{\ell,n}^{(m)}(\theta)
    \right],
    \label{eq:fendpointdisc}
\end{equation}
where the unnormalized target-state weights are
\begin{equation}
    a_{\ell,n}^{(m)}(\theta)
    =
    \exp \left[
        -u_\ell(\mathbf{x}_n, V_n, m; \theta)
        -
        \log Z_{\mathrm{mix},\ell}(n)
    \right].
    \label{eq:am}
\end{equation}
Because the prefactor $1/M_\ell$ is the same for every state of the same leg,
the implementation evaluates
\begin{equation}
    \widetilde{F}_{\ell,m}(\theta)
    =
    -
    \log
    \sum_{n=1}^{M_\ell}
    a_{\ell,n}^{(m)}(\theta),
    \label{eq:fendpointimpl}
\end{equation}
which differs from Equation \eqref{eq:fendpointdisc} only by the additive
constant $\log M_\ell$ and therefore produces the same endpoint free-energy
differences.

The normalized MBAR weights used to reweight observables to the target state $m$
are
\begin{equation}
    \omega_{\ell,n}^{(m)}(\theta)
    =
    \frac{
        a_{\ell,n}^{(m)}(\theta)
    }{
        \sum_{j=1}^{M_\ell}
        a_{\ell,j}^{(m)}(\theta)
    }.
    \label{eq:omegam}
\end{equation}
These are the weights implicitly formed by the current endpoint-gradient
routine in the code.

The same machinery can be applied to every lambda state in the schedule to
reconstruct a full replayed free-energy profile. Stage B readiness uses this
full-profile reconstruction to compare MBAR-replayed profiles against the frozen
AWH profile. Optimization, however, only needs the endpoint states.

\section{Endpoint Free-Energy Gradients}

The endpoint gradient estimator follows directly from the MBAR expression above.
Differentiating Equation \eqref{eq:fendpointimpl} with respect to $\theta$ gives
\begin{align}
    \nabla_\theta \widetilde{F}_{\ell,m}(\theta)
    &=
    -\frac{1}{
        \sum_{n=1}^{M_\ell} a_{\ell,n}^{(m)}(\theta)
    }
    \sum_{n=1}^{M_\ell}
    \nabla_\theta a_{\ell,n}^{(m)}(\theta)
    \nonumber \\
    &=
    -\frac{1}{
        \sum_{n=1}^{M_\ell} a_{\ell,n}^{(m)}(\theta)
    }
    \sum_{n=1}^{M_\ell}
    \left[
        -a_{\ell,n}^{(m)}(\theta)
        \nabla_\theta
        u_\ell(\mathbf{x}_n, V_n, m; \theta)
    \right]
    \nonumber \\
    &=
    \sum_{n=1}^{M_\ell}
    \omega_{\ell,n}^{(m)}(\theta)
    \nabla_\theta
    u_\ell(\mathbf{x}_n, V_n, m; \theta).
    \label{eq:fgradred}
\end{align}
Since the pressure-volume term in Equation \eqref{eq:reducedpotential} does not
depend on the force-field parameters,
\begin{equation}
    \nabla_\theta
    u_\ell(\mathbf{x}_n, V_n, m; \theta)
    =
    \beta_\ell
    \nabla_\theta
    U_\ell(\mathbf{x}_n, m; \theta),
    \label{eq:gradreduced}
\end{equation}
and therefore
\begin{equation}
    \nabla_\theta \widetilde{F}_{\ell,m}(\theta)
    =
    \sum_{n=1}^{M_\ell}
    \omega_{\ell,n}^{(m)}(\theta)
    \,
    \beta_\ell
    \nabla_\theta
    U_\ell(\mathbf{x}_n, m; \theta).
    \label{eq:fgradtheta}
\end{equation}
This is the precise sense in which the MBAR weights are used to compute endpoint
gradients in the implementation: every stored frame is replayed at the endpoint
lambda state of interest, and the corresponding force-field gradient at that
endpoint is averaged with the normalized target-state weights
$\omega_{\ell,n}^{(m)}(\theta)$.

The endpoint free energy of a leg is the coupled-minus-decoupled difference
\begin{equation}
    \Delta G_\ell(\theta)
    =
    \widetilde{F}_{\ell,\mathrm{coupled}}(\theta)
    -
    \widetilde{F}_{\ell,\mathrm{decoupled}}(\theta),
    \label{eq:legdg}
\end{equation}
and its gradient is
\begin{equation}
    \nabla_\theta \Delta G_\ell(\theta)
    =
    \nabla_\theta \widetilde{F}_{\ell,\mathrm{coupled}}(\theta)
    -
    \nabla_\theta \widetilde{F}_{\ell,\mathrm{decoupled}}(\theta).
    \label{eq:legdggrad}
\end{equation}
For a coefficient-weighted thermodynamic cycle,
\begin{equation}
    \Delta G_{\mathrm{cyc}}(\theta)
    =
    \Delta G_{\mathrm{std}}
    +
    \sum_\ell c_\ell \Delta G_\ell(\theta),
    \label{eq:cycle}
\end{equation}
and the residual against experiment is
\begin{equation}
    \mathcal{E}(\theta)
    =
    \Delta G_{\mathrm{cyc}}(\theta)
    -
    \widetilde{\Delta G}_{\mathrm{exp}}.
    \label{eq:residual}
\end{equation}
The code applies a Huber derivative to this residual, so the optimization
gradient is
\begin{equation}
    \nabla_\theta \mathcal{L}
    =
    \frac{d \mathcal{L}}{d \mathcal{E}}
    \nabla_\theta \Delta G_{\mathrm{cyc}}(\theta).
    \label{eq:lossgradtheta}
\end{equation}

\section{Bounded Coordinates}

The optimizer does not update unconstrained physical parameters directly. Each
trainable parameter is represented by a latent coordinate $\phi_i$ and mapped
into a bounded interval
$[\theta_{i,\min}, \theta_{i,\max}]$ through a shifted sigmoid transform,
\begin{equation}
    \theta_i(\phi_i)
    =
    \theta_{i,\min}
    +
    \left(
        \theta_{i,\max} - \theta_{i,\min}
    \right)
    \sigma(k \phi_i),
    \label{eq:phitheta}
\end{equation}
where $k$ is a configurable slope and $\sigma$ is the logistic function. The
chain rule then gives
\begin{equation}
    \nabla_{\phi_i} U
    =
    \frac{d \theta_i}{d \phi_i}
    \nabla_{\theta_i} U,
    \label{eq:chainphi}
\end{equation}
so all optimization-space gradients and Fisher matrices are assembled in latent
coordinates.

\section{From Euclidean to Natural Gradient Descent}

At this point the reason for introducing a metric tensor can be stated
explicitly. If one ignored the geometry of the probability distributions and
performed ordinary Euclidean gradient descent in latent parameter space, the step
would be
\begin{equation}
    \Delta \phi_{\mathrm{Euc}}
    =
    -
    \eta
    \nabla_\phi \mathcal{L}.
    \label{eq:eucstep}
\end{equation}
This update treats all coordinate directions in $\phi$-space as geometrically
equivalent. That assumption is poor for force-field optimization. Different
parameters perturb the sampled distribution with very different strength: a small
change in one parameter can produce a large statistical distortion, while a
numerically larger change in another parameter can have only a modest effect. In
other words, the relevant geometry is not the Euclidean geometry of the
coordinates themselves, but the geometry of the manifold of probability
distributions induced by those coordinates.

The natural gradient is the steepest-descent direction on that statistical
manifold. If $G(\phi)$ is the metric tensor, then the corresponding
natural-gradient update is
\begin{equation}
    \Delta \phi_{\mathrm{nat}}
    =
    -
    \eta
    G(\phi)^{-1}
    \nabla_\phi \mathcal{L}.
    \label{eq:natstep}
\end{equation}
In our setting, the appropriate local metric is the Fisher information matrix.
This does not mean that the code explicitly integrates exact geodesics of the
manifold. Rather, it uses the Fisher matrix as a local quadratic description of
the manifold curvature, so that the update direction is the local
steepest-descent direction consistent with distribution space rather than raw
coordinate space.

This is exactly how the current implementation uses the Fisher matrix. The code
first builds a latent-space loss gradient, then preconditions it by the empirical
Fisher matrix to obtain a natural-gradient-like base step. That step is then
stabilized by diagonal normalization, eigentruncation, KL-target scaling, and a
line search. The KL derivation below explains why this is the correct local
geometry to use.

\section{Why Fisher is the Metric Tensor}

The role of the Fisher information matrix is not merely heuristic. It appears as
the first nonzero local approximation to the Kullback-Leibler divergence between
nearby candidate ensembles. This is precisely why the current optimizer uses it
as a metric tensor: it measures how large a parameter step is in the space of
probability distributions, not in the raw coordinate system of the parameters.

For a fixed frozen reference Gibbs scheme $g_{\mathrm{ref},\ell}$ and current
latent parameter vector $\phi$, define the candidate extended-state density
\begin{equation}
    p_\phi(\mathbf{x}, V, i)
    =
    \frac{
        \exp \left[
            g_{\mathrm{ref},\ell}(i)
            -
            u_\ell(\mathbf{x}, V, i; \phi)
        \right]
    }{
        \mathcal{Z}_\ell(\phi)
    }.
    \label{eq:pphi}
\end{equation}
Here the dependence on $\phi$ enters only through the reduced potential; the
frozen AWH Gibbs log-weights are held fixed.

Now compare $p_\phi$ with a nearby candidate
$p_{\phi + \delta \phi}$. The KL divergence from the current ensemble to the
proposed ensemble is
\begin{equation}
    D_{\mathrm{KL}}
    \left(
        p_\phi
        \,\|\, 
        p_{\phi + \delta \phi}
    \right)
    =
    \mathbb{E}_\phi
    \left[
        \log p_\phi
        -
        \log p_{\phi + \delta \phi}
    \right].
    \label{eq:kl}
\end{equation}

Taylor-expand the log density of the proposed distribution around $\phi$:
\begin{equation}
    \log p_{\phi + \delta \phi}
    =
    \log p_\phi
    +
    \delta \phi^\top \mathbf{s}_\phi
    +
    \frac{1}{2}
    \delta \phi^\top
    H_\phi
    \delta \phi
    +
    O(\|\delta \phi\|^3),
    \label{eq:taylorlogp}
\end{equation}
where
\begin{equation}
    \mathbf{s}_\phi
    =
    \nabla_\phi \log p_\phi
    \label{eq:score}
\end{equation}
is the score vector and
\begin{equation}
    H_\phi
    =
    \nabla_\phi^2 \log p_\phi
    \label{eq:hessian}
\end{equation}
is the Hessian of the log density.

Substituting Equation \eqref{eq:taylorlogp} into
Equation \eqref{eq:kl} yields
\begin{equation}
    D_{\mathrm{KL}}
    \left(
        p_\phi
        \,\|\, 
        p_{\phi + \delta \phi}
    \right)
    =
    -
    \delta \phi^\top
    \mathbb{E}_\phi[\mathbf{s}_\phi]
    -
    \frac{1}{2}
    \delta \phi^\top
    \mathbb{E}_\phi[H_\phi]
    \delta \phi
    +
    O(\|\delta \phi\|^3).
    \label{eq:kltaylor}
\end{equation}

The score expectation vanishes:
\begin{equation}
    \mathbb{E}_\phi[\mathbf{s}_\phi]
    =
    \int p_\phi
    \nabla_\phi \log p_\phi
    =
    \int \nabla_\phi p_\phi
    =
    \nabla_\phi \int p_\phi
    =
    \nabla_\phi 1
    =
    0.
    \label{eq:scorezero}
\end{equation}
Therefore the linear term disappears and the \emph{first nonzero} contribution to
the local KL divergence is second order:
\begin{equation}
    D_{\mathrm{KL}}
    \left(
        p_\phi
        \,\|\, 
        p_{\phi + \delta \phi}
    \right)
    =
    \frac{1}{2}
    \delta \phi^\top
    \mathcal{I}(\phi)
    \delta \phi
    +
    O(\|\delta \phi\|^3),
    \label{eq:klmetric}
\end{equation}
where
\begin{equation}
    \mathcal{I}(\phi)
    =
    -
    \mathbb{E}_\phi[H_\phi]
    \label{eq:fimhessian}
\end{equation}
is the Fisher information matrix. Equation \eqref{eq:klmetric} is the reason the
Fisher matrix is the natural metric tensor: its quadratic form is the local
statistical distance induced by KL divergence.

\section{From the KL Expansion to a Covariance of Energy Gradients}

Under the usual regularity assumptions,
\begin{equation}
    \mathcal{I}(\phi)
    =
    \mathbb{E}_\phi
    \left[
        \mathbf{s}_\phi
        \mathbf{s}_\phi^\top
    \right].
    \label{eq:fimscore}
\end{equation}
To turn this abstract definition into the object actually used by the code, we
derive the score explicitly for the frozen Gibbs family in
Equation \eqref{eq:pphi}.

Taking the logarithm of Equation \eqref{eq:pphi} gives
\begin{equation}
    \log p_\phi(\mathbf{x}, V, i)
    =
    g_{\mathrm{ref},\ell}(i)
    -
    u_\ell(\mathbf{x}, V, i; \phi)
    -
    \log \mathcal{Z}_\ell(\phi).
    \label{eq:logpphi}
\end{equation}
Differentiating with respect to $\phi$ yields
\begin{equation}
    \mathbf{s}_\phi(\mathbf{x}, V, i)
    =
    -
    \nabla_\phi u_\ell(\mathbf{x}, V, i; \phi)
    -
    \nabla_\phi \log \mathcal{Z}_\ell(\phi).
    \label{eq:score1}
\end{equation}
Now
\begin{align}
    \nabla_\phi \log \mathcal{Z}_\ell(\phi)
    &=
    \frac{1}{\mathcal{Z}_\ell(\phi)}
    \nabla_\phi \mathcal{Z}_\ell(\phi)
    \nonumber \\
    &=
    \frac{1}{\mathcal{Z}_\ell(\phi)}
    \sum_i \int dV \int d\mathbf{x}
    \left[
        -
        \nabla_\phi u_\ell(\mathbf{x}, V, i; \phi)
    \right]
    \nonumber \\
    &\qquad \times
    \exp \left[
        g_{\mathrm{ref},\ell}(i)
        -
        u_\ell(\mathbf{x}, V, i; \phi)
    \right]
    \nonumber \\
    &=
    -
    \mathbb{E}_\phi
    \left[
        \nabla_\phi u_\ell(\mathbf{x}, V, i; \phi)
    \right].
    \label{eq:gradlogz}
\end{align}
Substituting Equation \eqref{eq:gradlogz} into Equation \eqref{eq:score1}
gives
\begin{equation}
    \mathbf{s}_\phi(\mathbf{x}, V, i)
    =
    -
    \nabla_\phi u_\ell(\mathbf{x}, V, i; \phi)
    +
    \mathbb{E}_\phi
    \left[
        \nabla_\phi u_\ell(\mathbf{x}, V, i; \phi)
    \right].
    \label{eq:score2}
\end{equation}
This is already a centered quantity, so the Fisher matrix is the covariance of
reduced-potential gradients:
\begin{equation}
    \mathcal{I}(\phi)
    =
    \mathrm{Cov}_\phi
    \left(
        \nabla_\phi u_\ell(\mathbf{x}, V, i; \phi)
    \right).
    \label{eq:fimcovu}
\end{equation}

Finally, because
$u_\ell(\mathbf{x}, V, i; \phi) = \beta_\ell U_\ell(\mathbf{x}, i; \phi) + \beta_\ell P_\ell V$
and the $pV$ term does not depend on $\phi$,
\begin{equation}
    \nabla_\phi u_\ell(\mathbf{x}, V, i; \phi)
    =
    \beta_\ell
    \nabla_\phi U_\ell(\mathbf{x}, i; \phi),
    \label{eq:gradubeta}
\end{equation}
which leads to the form used conceptually by the implementation:
\begin{equation}
    \mathcal{I}(\phi)
    =
    \mathrm{Cov}_\phi
    \left(
        \beta_\ell
        \nabla_\phi U_\ell(\mathbf{x}, i; \phi)
    \right).
    \label{eq:fimcovgradu}
\end{equation}
Thus the Fisher matrix is not an arbitrary preconditioner. It is the covariance
of the reduced-potential score vectors and the local metric tensor induced by KL
divergence.

\section{Empirical Reweighting, ESS, and the Practical Fisher Estimate}

The formulas above are expectations under the current candidate distribution
$p_\phi$, but the stored frames come from the frozen reference distribution
$p_{\mathrm{ref}}$. The code therefore reweights the reference frames into the
current candidate ensemble.

For the full extended state $(\mathbf{x}, V, \lambda)$,
\begin{equation}
    \frac{
        p_\phi(\mathbf{x}, V, \lambda)
    }{
        p_{\mathrm{ref}}(\mathbf{x}, V, \lambda)
    }
    =
    \frac{
        \mathcal{Z}_{\mathrm{ref}}
    }{
        \mathcal{Z}_\phi
    }
    \exp \left[
        -
        u_\ell(\mathbf{x}, V, \lambda; \phi)
        +
        u_{\mathrm{ref},\ell}(\mathbf{x}, V, \lambda)
    \right].
    \label{eq:ratio}
\end{equation}
Because the frozen Gibbs factor
$g_{\mathrm{ref},\ell}(i_n)$ appears in both numerator and denominator, it
cancels exactly. For a stored frame $n$ with visited lambda state $i_n$,
the unnormalized reweighting factor used by the optimizer is therefore
\begin{equation}
    W_{\ell,n}(\phi)
    \propto
    \exp \left[
        -
        u_\ell(\mathbf{x}_n, V_n, i_n; \phi)
        +
        u_{\mathrm{ref},\ell}(\mathbf{x}_n, V_n, i_n)
    \right].
    \label{eq:wactive}
\end{equation}
For NPT legs, the $pV$ term also cancels because the same pressure is used in the
reference and candidate ensembles. The normalized weights are
\begin{equation}
    \overline{W}_{\ell,n}(\phi)
    =
    \frac{
        W_{\ell,n}(\phi)
    }{
        \sum_{j=1}^{M_\ell} W_{\ell,j}(\phi)
    },
    \label{eq:wnorm}
\end{equation}
and the effective sample size is
\begin{equation}
    \mathrm{ESS}_\ell(\phi)
    =
    \frac{1}{
        \sum_{n=1}^{M_\ell}
        \overline{W}_{\ell,n}(\phi)^2
    }.
    \label{eq:ess}
\end{equation}

The empirical score vector for frame $n$ is taken at the physically visited
lambda state,
\begin{equation}
    \mathbf{s}_{\ell,n}(\phi)
    =
    \beta_\ell
    \nabla_\phi
    U_\ell(\mathbf{x}_n, i_n; \phi),
    \label{eq:scoreemp}
\end{equation}
and the weighted empirical Fisher matrix is
\begin{equation}
    \widehat{\mathcal{I}}_\ell(\phi)
    =
    \sum_{n=1}^{M_\ell}
    \overline{W}_{\ell,n}(\phi)
    \left(
        \mathbf{s}_{\ell,n}(\phi)
        -
        \overline{\mathbf{s}}_\ell(\phi)
    \right)
    \left(
        \mathbf{s}_{\ell,n}(\phi)
        -
        \overline{\mathbf{s}}_\ell(\phi)
    \right)^\top,
    \label{eq:fimemp}
\end{equation}
with
\begin{equation}
    \overline{\mathbf{s}}_\ell(\phi)
    =
    \sum_{n=1}^{M_\ell}
    \overline{W}_{\ell,n}(\phi)
    \mathbf{s}_{\ell,n}(\phi).
    \label{eq:meanemp}
\end{equation}
The full cycle Fisher matrix is the sum of the leg contributions in the shared
latent parameter space.

Because the first nonzero KL term is
$\frac{1}{2} \delta \phi^\top \mathcal{I}(\phi) \delta \phi$, the optimizer uses
the empirical Fisher matrix to define a trust region,
\begin{equation}
    D_{\mathrm{KL}}
    \approx
    \frac{1}{2}
    \Delta \phi^\top
    \widehat{\mathcal{I}}(\phi)
    \Delta \phi.
    \label{eq:kltrust}
\end{equation}
In the code, the empirical Fisher matrix is therefore not a decorative
diagnostic. It is the operator that turns the Euclidean loss gradient into a
natural-gradient-like base step, after which the optimizer applies diagonal
normalization, eigentruncation, and KL-target scaling before the line search.
The full logic is:
\begin{enumerate}
    \item replay the frozen reference frames to obtain endpoint free energies and
    endpoint gradients,
    \item assemble the latent-space loss gradient from those endpoint
    contributions,
    \item estimate the empirical Fisher matrix from reweighted score vectors at
    the visited lambda states,
    \item use that Fisher matrix as the local metric tensor of the statistical
    manifold,
    \item bound the proposed update through the quadratic KL model in
    Equation \eqref{eq:kltrust}.
\end{enumerate}
This is why a KL target, rather than a raw Euclidean step length, is the natural
quantity to bound. A step that is small in Euclidean coordinates can still be
large in distribution space if it distorts the ensemble strongly, while a step
that is large in a weakly constrained coordinate may still be statistically
safe. The Fisher metric captures that distinction directly.
